\selectlanguage{english}\Traditional
\Language{zh-Hant}{繁體中文}{使用指南}
\firsttopic{啟動、休眠與返回筆記}
Inkline 在 reMarkable 2 上運行本機 shell。使用 Type Folio 或 USB 鍵盤，可以在平板上
工作，也可以通過 SSH 或 Goblin Mosh 連線其他計算機。文字和圖像以適合電子紙的灰度顯示。

\begin{notice}
按 \key{Opt+RightAlt+T}，即可直接在平板上打開 Inkline。\par
\key{電源鍵：}按下再松開即可休眠，再按一次喚醒。已打開的 shell 和編輯器緩衝區保留在 RAM 中。
\end{notice}
休眠會暫停設備，不會退出 Inkline，也不會把設定寫入快閃記憶體。喚醒後，網路連線可能需要重新建立。
用於喚醒的按鍵不會再次觸發休眠。更新後，請退出並重新打開 Inkline 一次，以啟用新的電源鍵處理。

要返回筆記，點選 \key{Quit} 或按 \key{Ctrl+Shift+Q}，然後確認。
輸入 \code{exit} 只關閉當前終端；關閉最後一個終端後返回筆記。
底欄包含 Escape、電池狀態、Quit 和 Settings。
啟動時的簡短說明和小 Goblin 標誌可用 \code{clear} 清除。

如果快捷鍵啟動的應用佔用了屏幕，按 \key{Opt+RightAlt+T} 返回 Inkline。
緊急恢復時，\key{按住 Opt+RightAlt+Backspace 兩秒}。
此操作會關閉所有 Inkline 終端，丟失這些會話中尚未儲存的工作。
設備重新啟動後仍預設進入筆記界面。

RightAlt 指空白鍵右側的 Alt/Opt 鍵。按住它與獨立的 Opt 鍵，再按 T；不要按 Ctrl。單獨使用 Super 的組合仍可自訂。

\ProjectLinks

\topic{按鍵與終端槽位}
Type Folio 在 Ctrl 與 Alt 之間有一個獨立的 \key{Opt} 鍵，另有一個
\key{右 Alt/Opt} 鍵。兩者用途不同。

\begin{keytable}
\key{按鍵} & \key{操作} \\ \midrule
獨立 Opt + 1--0 & F1--F10 \\
右 Alt/Opt + 1--9 & 選擇終端槽位1--9 \\
右 Alt/Opt + Left / Right & 上一個／下一個槽位 \\
右 Alt/Opt + Tab & Escape \\
右 Alt/Opt + Up / Down & PageUp / PageDown \\
右 Alt/Opt + Backspace & 確認退出所有終端 \\
右 Alt/Opt + Space & Settings \\
左 Alt + Space & Unicode 鍵盤 \\
右 Alt/Opt + C / V & 複製選區／貼上 \\
\end{keytable}

Ctrl+Shift+B 會原樣傳給終端程式，包括 Emacs。

在\key{美式 Type Folio} 上，右 Alt/Opt+\key{0} 輸入 \code{+}；
右 Alt/Opt+\key{緊挨 Backspace 左側的減號鍵}輸入 \code{=}。
關閉輸入法後，Shift+6 直接輸入插入符號：\code{c^2} 保持為 \code{c^2}。
沒有用於縮放字體的鍵盤組合鍵。

USB 鍵盤可直接使用普通功能鍵。獨立 Opt 的快捷操作在提供 Meta 的 USB 鍵盤上使用 Meta 修飾鍵。

最多有\key{九個獨立槽位}。選擇空槽位會打開一個 shell。指示器顯示當前打開的終端數，而非容量。
如果只有槽位1和9打開，在9中顯示 \key{Terminal 2 / 2 · slot 9}。
關閉 shell 會釋放槽位，不會重新編號其他快捷鍵。

每個終端都有自己的 shell、工作目錄、輸入組合狀態、圖像和500行回捲歷史。
顯示設定和輸入法設定對所有槽位生效。

\topic{設定與顯示}
點選底欄最右側的 \key{Settings}，或按 \key{右 Alt/Opt+Space}；底欄隱藏時也可使用。
\key{左 Alt+Space} 開啟 Unicode 鍵盤。實心選項表示當前狀態，虛線邊框表示鍵盤焦點。
Tab 移動焦點，方向鍵和 Enter 更改選擇。

\key{Caps Lock：}可選擇 Control（預設）或 Caps Lock。
\key{字體大小：}雙指捏合，或使用設定中的減號／加號。
緩慢移動可每次調整一個像素，範圍為\key{6至48像素}。

\key{Text darkness：}50\% 為原始渲染，數值越高，字形邊緣越深。
\key{Minimum contrast} 在應用選擇相近顏色時拉開前景與背景的亮度差，預設35\%。

\begin{choices}
\key{電子紙模式} & \key{取捨} \\ \midrule
Fast & 合併等待最少，使用動畫波形，殘影較多 \\
Balanced & UI 波形，按32 ms 合併輸出 \\
Crisp（預設） & 灰度更清晰，使用內容波形，更新較慢 \\
Mono & 快速黑白顯示，捨棄圖像中的灰色層次 \\
Saver & 按120 ms 合併，減少連續輸出的刷新次數 \\
\end{choices}

只重繪發生變化的文字行，但面板本身仍需要刷新時間。
Saver 會減少刷新請求，省電效果尚未測量。升級會保留明確儲存的模式，不會強制覆蓋為新預設值。

\begin{notice}
所有設定在使用期間保留在 \key{RAM}，正常退出 Inkline 時一起儲存。
抬起手指、關閉設定和休眠都不會寫入快閃記憶體。沒有實際更改的會話不會寫入設定。
崩潰、強制終止或斷電可能丟失尚未儲存的設定更改。
\end{notice}

底欄每分鐘從唯讀的 Linux 電源介面讀取一次電池狀態。\code{inkline-battery}
顯示電量和充電狀態；\code{inkline-battery --percentage} 或
\code{~/inkline battery} 顯示 Inkline 無色 Bash 提示符所用的簡短數值。

\topic{Unicode 與輸入法}
按 \key{左 Alt+Space} 打開 Unicode 鍵盤。點選字元即可輸入，鍵盤保持打開，方便連續輸入。
點選 \key{Done} 或按 Escape 關閉。

分類包括字母、組合標記、數字、標點、數學／貨幣、符號、空格、格式字元和私用區。
上下拖動滾動；也可以使用方向鍵、PageUp/PageDown、Home/End 或滑鼠滾輪。
\key{Tab} 切換分類，Enter 輸入選中的字元。

輸入十六進制 Unicode 值（例如 \code{03bb}）後按 Enter，可以輸入 λ。
Backspace 編輯數值。組合標記旁的虛線圓僅用於展示，實際只輸入標記。
空格和格式字元帶有可見標籤。字元表使用 Qt 的 Unicode 數據，不包括控制字元、
未分配碼位、代理碼位和非字元。部分字元需要未安裝的字體，私用區需要匹配的字體。

在鍵盤中按 \key{F2} 或點選 \key{Settings}，然後選擇 \key{Input method}：
\begin{choices}
Off — direct keyboard & 普通美式鍵盤直接輸入 \\
Romaji & 從羅馬字輸入日語平假名 \\
Pinyin & 拼音輸入中文 \\
Zhuyin & 使用基本注音鍵位輸入中文 \\
Wubi 86 & 使用五筆86字形編碼輸入中文 \\
US-International & 常用拉丁字母重音 \\
\end{choices}

在候選欄選擇結果。要恢復普通美式輸入，在同一菜單選擇
\key{Off — direct keyboard}。此選擇對所有終端生效，並在 Inkline 退出時儲存。

\topic{觸控、手寫筆與剪貼簿}
\key{雙指上下滑動：}滾動歷史。向下拖動查看較早輸出，向上返回提示符方向。
每個終端除當前屏幕外，還在 RAM 中保留500個物理行。

\key{雙指左右滑動：}每次手勢切換一個終端槽位。再次切換前請抬起手指。
雙指張開或合攏可調整字體大小。

\key{單指：}拖動選擇文字。右 Alt/Opt+C 複製，右 Alt/Opt+V 貼上。
所有終端共享一個最大128 KiB 的 RAM 剪貼簿。

\key{手寫筆：}啟用終端滑鼠報告的應用會收到對應的滑鼠左鍵按下、移動與松開事件，
使用應用協商的協議，包括 SGR。其他情況下拖動用於選擇文字。
按住 \key{Shift}，可以在啟用滑鼠報告的應用中強制進行本機選擇。

\key{OSC 52} 允許終端程式（包括經 SSH 運行的遠端程式）讀取或替換 Inkline 的 RAM 剪貼簿。
它與筆記應用的剪貼簿不同。右 Alt/Opt+X 會複製，並提示在編輯器內部執行刪除；
OSC 52 無法刪除編輯器里的文字。

\key{圖像：}kitty 的內聯與共享記憶體傳輸保留在 RAM 中；檔案與臨時檔案傳輸禁用。
顯式請求時支持 sixel，但尚未自動聲明支持。kitty 佔位符、動畫播放和部分 sixel 模式尚未完成。

\code{clear} 和全屏擦除會同時清除可見圖像與文字。記憶體不足時會回收屏幕外的圖像。
沒有磁盤圖像快取，但大圖仍會佔用設備有限的 RAM。

\topic{USB 鍵盤與打字機}
開啟\key{Settings → USB（第 3 頁）}。\key{Network} 恢復 USB 乙太網路；
\key{Send keyboard} 讓平板成為另一台電腦的鍵盤；\key{Receive keyboard}
透過有供電的 OTG 轉接器接收有線鍵盤。Type Folio 在各模式下均可使用。
輸出鍵盤模式需要先安裝網站提供的 Python 工具，無需額外的 Python 模組。

選擇 Send keyboard，以 USB-C 連接電腦，設定下列接收端配置，然後開啟
\key{Open typewriter}。Send keyboard 已啟用時會自動開始即時傳送。
暫停時編輯器以黑底清楚顯示 \key{PAUSED · USB OUTPUT OFF}；即時傳送使用另一種狀態標誌。
\key{Escape} 或\key{Stop} 暫停並釋放按鍵。
離開打字機、切換終端機、視窗失焦、開啟原生應用程式或休眠時也會暫停。
斷線重連後不會自動重送。

\begin{choices}
US / ASCII & 美式配置，關閉 Caps Lock。拒絕非 ASCII 文字。 \\
Mac & 在鍵盤輸入來源中選擇 Unicode Hex Input，關閉 Caps Lock。
只支援到 U+FFFF，拒絕補充平面字元。 \\
Linux & 美式配置，關閉 Caps Lock。接收應用程式／輸入法必須支援 Ctrl+Shift+U、
十六進位數字、Enter；並非每個應用程式都支援。 \\
Windows & 在 PC 上安裝 WinCompose，將 Compose 設為 Right Alt，啟用 Unicode 輸入。
使用美式配置並關閉 Caps Lock。傳送 Compose、u、六位十六進位數字、Enter。 \\
\end{choices}
從 \url{https://wincompose.info/} 取得 WinCompose。自訂規則不可覆寫此 Unicode 序列。
HID 沒有通用的 Unicode 文字通道；結果取決於接收端設定與字型。先用短文字測試。

\key{Input method} 與終端機共用輸入引擎：Off／直接美式、Romaji、Pinyin、Zhuyin、
Wubi 86 和 US-International。只有已確認字元才會送出。更換輸入法後需點 Start 繼續。
選擇\key{Off} 返回普通美式輸入。\key{Unicode} 或左 Alt+Space 開啟的字元面板也向 USB 接收端輸入。

Typewriter 是本機編輯器，Inkline 執行期間文件只保存在 RAM。一般輸入會編輯本機文件；
開啟即時傳送後，每個已確認字元也會傳送到連線的電腦。複製、剪下、貼上、選取和雙指捲動
都作用於本機編輯器。\key{Files \& send} 可用任何檔名載入或儲存 UTF-8，
以每秒 5--80 個字元傳送全文，並可停止傳送。編輯時不會寫入儲存空間；只有 Save 或
Inkline 正常結束時才會寫入。無標題文件在結束時使用隱藏復原名
\path{~/.inkline-typewriter-draft}。\key{Exit typewriter} 返回 USB 設定。

\topic{傳送檔案與恢復網路}
安裝器把 \code{inkline-usb} 和 \code{keyboard-send} 放入
\path{~/.local/bin}。\code{~/inkline usb}、\code{~/inkline type} 和
\code{inkline-type} 保留為相容別名。Outpost 的獨立工具保持原樣。
\UsbExamples
UTF-8 檔案\key{沒有固定大小上限}。傳送器分塊讀取，不在 RAM 中儲存整份文件，
也不建立中轉檔案。可回讀的一般檔案會先完整驗證再重新讀取傳送；傳送時不要修改來源檔案。
管線按到達順序驗證，因此後續錯誤出現時前面的內容可能已經送出。
可選的 UTF-8 BOM 會被去除，CRLF／CR 轉為換行。Enter 和 Tab 會操作目前的應用程式。

\code{--cps}（預設 20）或 \code{--delay-ms} 設定速度；Unicode 需要更多按鍵報告。
\code{--start-delay} 留出切換輸入焦點的時間。Dry-run 只列印報告，不傳送 USB 按鍵。
Ctrl+C 或 \code{--stop} 取消並釋放按鍵。同時只能有一個傳送器。
互動佇列滿時會暫停；中斷後重試前先檢查接收文件。文字和暫存狀態不寫入快閃儲存。
接收端配置與其他設定一起在正常結束時儲存。

點\key{Network} 或執行 \code{inkline-usb network} 恢復乙太網路。
Inkline 結束、當機或緊急重新啟動時也會恢復網路。USB 模式不會儲存為開機偏好。
\UsbRecoveryExample
計時器獨立於 shell 執行。從網路模式切換出去應使用平板或 Wi-Fi SSH，不能使用 USB SSH。
切換失敗會嘗試恢復乙太網路。Inkline 不會搶占使用中的衛星 PPP 或未知 USB 裝置功能；
結束清理也遵循此規則。Opt+RightAlt+Backspace 緊急重新啟動仍然可用，
但會關閉所有終端機並遺失其中未儲存的工作。

\topic{用快捷鍵啟動程式}
在 Inkline shell 中註冊快捷鍵。參數按字面傳遞；如需 shell 語法，請顯式調用 shell。
快捷鍵程式通過 Bash 讀取 \code{~/.bashrc}。

Settings → Command shortcuts（第 2 頁）可為字母鍵指定命令。選擇 Terminal 或 Native e-paper app，再按 Save。Remove 需確認後才刪除；Reset draft 放棄編輯。T 已鎖定。只有 Save 和確認刪除會寫入儲存空間。命令列工具管理同一組快捷鍵。

所有全域快捷鍵現在使用 Opt+RightAlt，包括 T 和長按 Backspace 的緊急復原組合。Folio 使用獨立的 Opt 鍵；USB 鍵盤使用 Windows/Command（Super）鍵。一般 Ctrl+Alt 仍可用於 Emacs。升級後請結束並重新開啟 Inkline；安裝會保留已開啟的終端機。

\begin{shell}
~/inkline shortcut register e emacs
~/inkline shortcut register p python3
~/inkline shortcut list
~/inkline shortcut deregister p
\end{shell}

\key{Opt+RightAlt+E} 在空閒槽位打開 Emacs。
\code{~/inkline run PROGRAM ARG...} 可不註冊直接啟動。
重複註冊會替換綁定；取消註冊不影響正在運行的程式。
按鍵採用美式鍵盤的物理位置。標點請加引號，自定義程式請使用完整路徑，並先安裝程式。

對於自行繪制屏幕的原生 Qt／電子紙應用：
\begin{shell}
~/inkline shortcut register a --epaper /home/root/my-app
\end{shell}
一次運行一個原生應用，期間暫停 Inkline 及其 shell。啟動器設定 Qt 電子紙環境與 RAM 快取，
丟棄日誌，並在應用退出後恢復 Inkline。應用應使用 Qt 輸入，不要獨佔鍵盤設備。

\begin{notice}
\key{Opt+RightAlt+T 是固定快捷鍵，不能註冊或取消註冊。}
它會停止原生快捷鍵應用，並返回現有終端。
此組合不需要 Ctrl；Caps Lock 映射只在 Inkline 內部生效。
\end{notice}

\key{緊急恢復：}按住 Opt+RightAlt+Backspace 兩秒，停止快捷鍵應用及其子程序並重新啟動 Inkline。
所有終端會關閉，未儲存的工作會丟失。以 root 運行的應用可能通過獨佔輸入、修改服務或耗盡資源
使恢復失敗；最後的辦法是重新啟動設備。

\topic{安裝、升級與保留手冊}
此安裝包面向\key{運行韌體3.27的 reMarkable 2}，使用系統自帶 Qt 6.8 電子紙庫。
最新下載與校驗步驟見：

\InstallLink

安裝器檢查硬件、韌體、字體、Qt 啟動、RAM 臨時目錄和 swap 是否禁用。
版本儲存在 \path{/home/root/.local/share/inkline} 下，鍵盤啟動服務隨系統啟動。
更新不會關閉已有終端。\key{退出並重新打開 Inkline 才會使用新版本。}
舊版本保留用於回退，會繼續佔用儲存空間。

安裝包附帶這一個包含九種語言的 PDF。筆記界面正在運行且 USB 網頁接口已啟用時，
安裝器會通過該接口嘗試導入\key{My files}。成功導入後記錄校驗值，避免重複導入同一手冊。

如果自動導入不可用，請連線 USB、退出 Inkline，並在筆記界面的儲存設定中啟用
\key{USB web interface}，然後通過 SSH 運行：
\begin{shell}
~/inkline manual
\end{shell}
也可下載 PDF，通過 USB 瀏覽器界面或 reMarkable 桌面／移動應用導入。
導入器不會直接修改筆記數據庫，也不會中斷終端。
卸載 Inkline 後，已導入手冊仍留在 My files。

通過 SSH 返回筆記或卸載 Inkline：
\begin{shell}
~/inkline stop
~/inkline uninstall
\end{shell}
可選軟體（包括推薦的 LaTeX 環境）的安裝說明保留在網站上。

\topic{工作原理、儲存與 RAM}
Qt 將鍵盤事件交給共享映射器，處理快捷鍵與 Caps Lock。
libghostty-vt 編碼終端按鍵，每個槽位通過獨立偽終端（PTY）連線程式。
shell 獲得 \code{HOME=/home/root}，交互式 Bash 讀取 \code{~/.bashrc}。

程式輸出經過終端解析器與 sixel 解碼器。libghostty 管理字元單元、光標狀態、備用屏幕、
回捲歷史和 kitty 圖像。渲染器保留灰度圖，只重繪變化的行；系統 Qt 電子紙後端負責刷新面板。
獨立的 evdev 守護程序處理全局快捷鍵和電源鍵，systemd 負責休眠與喚醒。

每個終端的核心分配上限為64 MiB；kitty 圖像每個屏幕16 MiB；保留的 sixel 圖像16 MiB。
渲染器和應用記憶體另計。空槽位不分配終端。九個大型應用可能超過可用 RAM。

臨時圖像、快取和會話狀態使用 RAM。已安裝軟體、儲存的文件、退出時儲存的設定以及應用寫入的
檔案使用持久儲存。Inkline 不會禁止任意程式寫檔案。

可選 Python 包是標準 Python。為減少快取寫入，請在設備的 \code{~/.bashrc} 中添加：
\begin{shell}
export PYTHONDONTWRITEBYTECODE=1
export PIP_NO_CACHE_DIR=1
\end{shell}
運行 \code{source ~/.bashrc} 或打開新 shell。
使用 \code{python3 -m pip install --no-compile PACKAGE} 可避免 pip 在安裝時單獨進行的位元組碼編譯。
Python 隔離模式會忽略 Python 環境變量。詳細內容見網站的 Python 指南。

Inkline 採用 GPL-3.0-or-later 許可。安裝包包含匹配的原始碼歸檔與依賴許可說明。
測試記錄區分自動檢查與仍需實機確認的行為。
